\clearpage
\chapter*{Đính chính cho Báo cáo số~212}
\addcontentsline{toc}{chapter}{Đính chính cho Báo cáo số~212}
\setcounter{footnote}{2}
\sid{e212}
\markboth{Đính chính 212}{Đính chính 212}
\src{301}

\noindent\textbf{5. TDTE III: Các tiền lược đồ thương.}

\begin{center}
[Séminaire Bourbaki, tập 13, 1960/61, số 212, 20 trang.]
\end{center}

\noindent\textbf{Trang 15, khẳng định (i).} --- Việc $Y=X/R$ là tựa xạ ảnh
trên $S$ dường như, cho đến nay, chỉ được chứng minh trong trường hợp $R$
phát sinh từ một quan hệ tương đương.

\noindent\textbf{Trang 18 và 19.} --- Như chúng tôi chỉ ra ở cuối báo cáo kế
tiếp, phỏng đoán đang xét hoàn toàn sai. ``Sự kiện khẳng định'' nêu trong Nhận
xét 8.1 dường như đã được nhiều tác giả chứng minh đồng thời (NAGATA,
ROSENLICHT, GROTHENDIECK, \ldots).

Xin lưu ý rằng việc áp dụng trong TDTE V lý thuyết chuyển sang thương được
phát triển ở đây vào việc xây dựng các lược đồ Picard cũng có thể được thay
thế bằng việc sử dụng thích hợp các lược đồ Hilbert (xem SHMT). Như đã nói
trong §~8, thiếu sót quan trọng nhất của lý thuyết được trình bày

\src{302}
ở đây là chưa có một tiêu chuẩn tồn tại cho các thương theo một quan hệ tương
đương không proper, chẳng hạn các quan hệ tương đương phát sinh từ một số tác
động của nhóm xạ ảnh. Một định lý quan trọng theo hướng này đã được MUMFORD
chứng minh\footnote{MUMFORD (David). --- An elementary theorem in geometric
invariant theory, \emph{Bull. Amer. Math. Soc.}, tập 67, 1961, tr. 483--486.}.
Về một sự tinh chỉnh kết quả của ông và nhiều ứng dụng vào lý thuyết, xem
SHMT.

\par
